\fancychapter{Introduction}
\label{chap:introduction}
% In total 3-4 pages

This thesis presents the design, simulation, implementation and experimental
validation of a custom \ac{mppt} charger for the photovoltaic system of
\ac{tsb}. The charger is intended to extract the maximum available power from
the team's solar panels under the rapidly changing irradiance of a marine
environment, while charging the low voltage battery of \ac{sg01} and providing
telemetry through a \ac{can} interface.

This chapter introduces the context of the work, the motivation for developing
an in-house \ac{mppt}, the project objectives, and the organization of the
remainder of this document.

\section{Motivation}
\label{sec:introduction:motivation}

With environmental pollution intensifying worldwide, the transition to cleaner energy sources such as solar power has become increasingly urgent.
In 1985, only 20.82\% of the energy produced in the world came from renewable energies. These numbers have been rising every year, and in 2024, have reached 31.92\%~\cite{ourworldindata_electricity_mix_2025}. In 2024, 6.9\% of the energy produced in the world come from solar panels, and in Portugal this number rises to 14.5\%, demonstrating the importance and the impact of solar energy~\cite{ourworldindata_electricity_mix_2025}. Solar energy still plays a minuscule role, and it is listed behind the other sources of energy in terms of its contribution to meeting the world's energy demand. However, solar energy is becoming more relevant, with the cost of solar panels dropping significantly in the last decade. 

Compared with other renewable sources, photovoltaic energy is advantageous owing to its availability, simplicity, low maintenance, environmental friendliness and reliability.
More recently, it is becoming more relevant in the automotive industry, with solar powered cars, boats, and robots~\cite{shahira2022_electrical_design_solar_boat}. The CO2 emissions of the automotive sector are one of the main contributors to global pollution. In 2018, 29\% of total CO2 emissions in the EU came from the transport sector, with 4\% coming from the maritime sector alone~\cite{icct_transport_eu_carbon_budget}.
These challenges have directly motivated the development of the \ac{tsb} project.

In 2015, \ac{tsb} was created with the goal of designing and building solar powered boats to compete in international competitions while also contributing to the innovation in yath industry. Since then, the project has grown, and several vessels have been built. It began with the construction of the first solar prototype, São Rafael 01, which had a lot of room for improvement and was followed by São Rafael 02 and 03. The team then decided to approach the hydrogen energy and autonomous driving, and therefore São Miguel 01 and 02 and São Pedro 01 were built. More recently, a vessel was built to combine all the technologies used before and was named \ac{sg01}.   

All these prototypes used solar energy to maximize their range and efficiency. In the first years the energy produced was small, and the whole system was built using commercial \ac{ots} components. However since the beginning of the \ac{tsb} project, the spirit of innovation and the desired to develop all the systems in house  has been present. In previous years, the team has built its own batteries, telemetry systems, flight controllers, fuel cell controllers, etc. Some other more complex systems were design during master thesis. Until today a \ac{bms}, dual motor system, Marine Propulsion System, a \ac{mppt} and more.

In 2020, the team started to build its own solar panels for São Rafael 03. To increase performance and the amount of information acquired by the solar panels' operation, in the same year, surge the first attempt to do a \ac{mppt}. This was the fist step to achieve a fully functional product. The result was a working \ac{mppt} that was never implemented on any vessel due to constantly burning components which was not reliable. Additionally, the project is outdated in terms of components and firmware, lacking configurability.


The \ac{mppt} charger is a fundamental subsystem in photovoltaic energy systems, responsible for maximizing the power extracted from solar panels by ensuring continuous operation at the \ac{mpp}. A \ac{mppt} charger is typically a power converter that continuously monitors the voltage and current produced by the photovoltaic modules and adjusts their electrical operating point accordingly. By dynamically regulating these operating parameters, the \ac{mppt} compensates for variations in environmental conditions such as solar irradiance and temperature, which cause the \ac{mpp} to shift over time. As a result, the photovoltaic system can operate with improved efficiency and deliver maximum available power under changing conditions~\cite{Comparative_study_MPPT_algorithms_2000}.

In the context \ac{tsb} project, a solar system is specially important to increase the autonomy of our vessels. In 2023/24 season we were able to consumption of \ac{sr03} from 1.5~kW to 300~w by using solar energy
Decreasing 80~\% of his consumption. With a costume made \ac{mppt}, we believe that we could reach even better performances wile also increasing the data acquired about the system performance that would serve to improve the system.

Now in 2025/26 season, the relevancy of the solar system decreased with our new vessel \ac{sg01}, since it uses around 20~KW of constant power. So the main power source will be hydrogen. However, solar energy will continue to be used in this vessel to power the auxiliary systems being a reliable green source that will keep the boat powered in all situations.

\section{Objectives} 

This project aims to design and implement a \ac{mppt} charger for the solar panels used in \ac{tsb} project. The \ac{mppt} charger will convert the energy produced by the solar panel as efficiently as possible with the use of a quality \acs{dc}-\acs{dc} converter and the implementation of \ac{mpp} tracking algorithms.

The main objectives of this project are:

\begin{itemize}[leftmargin=4em, itemsep=1pt, topsep=2pt]
    \item Study and understand the operation of solar panels and \ac{mppt} techniques;
    \item Choose the most suitable \acs{dc}-\acs{dc} topology for the system;
    \item Implement a fast response and accurate control algorithm;
    \item Design and simulate the hardware and software;
    \item Implement the system in hardware and develop a \ac{pcb};
    \item Test and validate the performance of the designed circuit;
    \item Ensure the safety and reliability of the \ac{mppt} charger for its integration in the \ac{tsb} project;
    \item Provide performance data of the solar panel and the \ac{mppt} through a \ac{can} communication interface;
    \item Develop a practical solution to change the prototype configuration.
\end{itemize}


By achieving these objectives, the project will contribute to a more efficient solar system, maximizing the energy received from the solar panels and increasing the available data for performance tracking of the \ac{mppt} and solar panel. This data will be used to later improve the energy efficiency of the \ac{mppt} and take conclusions about the manufacture quality of the solar panels build by \ac{tsb} project. A cheap, efficient and versatile \ac{mppt} will also mean future savings to the team, since it could be used in a variety of vessels.

In the context of \ac{tsb}, such complex projects with high scientific rigor, that involve simulations, \ac{pcb} design, laboratory tests and field tests are only possible in master thesis. So, for \ac{tsb}, the main objective is to increase the team knowledge about this subject by increasing the documentation about it. Also, due to a drastic event of another team, which burn all their \acp{mppt} during competition, the security of this system is the main concern.

\newpage


\section{Outline}

This document is organized into seven chapters:

\begin{itemize}[leftmargin=4em, itemsep=1pt, topsep=2pt]
    \item \textbf{Chapter~\ref{chap:introduction}} introduces the context of the work, presenting the motivation behind the development of a custom \acs{mppt} charger for the \acs{tsb} project. The objectives of the project are defined, and the structure of the document is outlined;

    \item \textbf{Chapter~\ref{chap:System}} describes the electrical architecture of \acs{sg01}, including the high and low voltage systems, the photovoltaic installation, and the batteries used in this work. The contribution of the solar system to vessel autonomy is analysed, and the time scales of environmental variations are estimated in order to define the required tracking speed of the \acs{mppt};

    \item \textbf{Chapter~\ref{sec:introduction:background}} reviews the State-of-the-Art of photovoltaic systems and \acs{mppt} technology. It covers the operating principles of solar panels, the most common \acs{mppt} algorithms, non isolated \acs{dc}-\acs{dc} converter topologies, lithium-ion battery charging techniques, and commercially available solutions;

    \item \textbf{Chapter~\ref{chap:imp}} presents the implementation of the proposed charger. The design of the synchronous boost converter, the selection of the processing unit and sensors, and the firmware architecture are described, including the charging state machine and communication interfaces;

    \item \textbf{Chapter~\ref{chap:sims}} reports the simulations used to validate the converter design and the \acs{mppt} algorithm. A photovoltaic cell model and the boost converter are simulated in LTspice, while an analytical power-loss model is developed in MATLAB to estimate the efficiency and the contribution of each component. Additionally a using MATLAB/simulink, the algorithm speed and accuracy was tested ;

    \item \textbf{Chapter~\ref{chap:lab_res}} presents the experimental validation of the prototype. Three hardware iterations are assembled and tested in the laboratory, covering sensing accuracy, converter efficiency, thermal behavior, \acs{mppt} tracking performance, and field tests;

    \item \textbf{Chapter~\ref{chap:Conclusion}} summarizes the main conclusions of this work and identifies directions for future development.
\end{itemize}
    